\scanpage{427}
\mspage{166}
Ainsi, le foncteur $\psi\varphi$ est bien homotope au foncteur identique. On a gagné !

Passons à $1')$, $2')$ pour $\mathcal C_{\mathbb N^{\unsure{o|0}}}$. L'assertion $1')$ se démontre exactement de la $=$ manière que 1) pour $\underline{\mathrm{Ch}}(X)$, en regardant $\mathcal C_{\mathbb N^o}$ comme fibrée sur $\deleted{\underline{\unsure{\mathrm{Tz}|\mathrm{Ty}|\mathrm{Iz}|\mathrm{Iy}}}}\mathbb N^o$ (qui a un objet final $0$), et en \deleted{\unsure{du}}utilisant le diagramme
\diagram{\begin{tikzcd}[ampersand replacement=\&,column sep=small]
X\arrow[r,"i'"]\arrow[d,equals] \& X_{\deleted{\underline{\mathrm{Tz}}}\mathbb N^o}=X\times\deleted{\underline{\mathrm{Tz}}}\mathbb N^o\arrow[rr,"j'"]\arrow[dr]\&\&\mathcal C_{\mathbb N^o}\arrow[dl]\\
\deleted{\unsure{=}}(X_{\deleted{\underline{\mathrm{Tz}}}\mathbb N^o})_{\unsure{e|0}}\&\&\deleted{\underline{\mathrm{Tz}}}\mathbb N^o\&{,}
\end{tikzcd}}
\deleted{\unsure{où}} où $j'$ est (cartésien et) une Hot-équivalence sur chaque fibre, d'après le lemme, donc une Hot équivalence \deleted{\unsure{part}} globalement, et $i'$ est une Hot-équivalence car $\mathbb N^o$ est asphérique (ayant un objet final). Donc $j'i'$ est asphérique qed.

Il reste à prouver que
\[
C_{\mathbb N^o}\mathrel{\unsure{\text{---}|\longrightarrow}}X\times X
\]
est propre et lisse. \deleted{Par \unsure{dualité}, \unsure{passage|passant} de $X$ à $X^{\unsure{o|'}}$, il suffit de prouver que c'est lisse\unsure{.|,}} Ici encore,
% archive sheet number: 426
